\item \label{itm:thermtec} Thermotectonic evolution \\
\begin{figure}[hbt]
    \centering
    \includegraphics[]{figures/tectonic_processes.eps}
    \caption{Thermal evolution of several idealized tectonic terranes.}
    \label{fig:thermtec}
\end{figure}

Both continental orogens and accretionary wedges are built structurally by
thrusting, yet the thermal evolution appears quite different.  Consider the
hypothetical thermal evolution of two tectonic environments presented in
Figure~\ref{fig:thermtec}.  A stable craton is provided as a reference. 

Describe the features and change in properties/sources that give rise to the
specific thermal evolution pattern related a thrust (Figure~\ref{fig:thermtec}).
\begin{enumerate}
    \item Why does heat flow initially decrease?
    \item Why does heat flow rebound?
    \item Why does heat flow reach a peak in an orogen, but not an accretionary
    wedge?
    \item Why does the long-term thermal evolution of a thrust take so long to
    approach that of a craton?
    \item Why should cratonic heat flow decrease slightly with time?
\end{enumerate}

As part of your discussion:
\begin{itemize}
    \item Draw the idealized geotherm immediately following the
    thermotectonic event.  Label this curve $t_0$.

    \item Using a different colored pen/pencil, draw the equilibrium
    geotherm.  Label this curve $t_\infty$.

    \item Using a third color, draw three hypothetical cooling at a few
    instances in time.  Label these curves $t_1$, $t_2$, and $t_3$.

    \item Your result will invariably look different than the model in
    Figure~\ref{fig:thermtec}.  What additional processes or changes to the
    lithosphere occur during a thrusting that may contribute to this
    difference?  Why?
\end{itemize}

{\bf Answer:} \\
This question draws on a fair amount of conceptual geologic and petrologic
knowledge and requires synthesizing those concepts with geophysical
concepts.  There is a fairly wide latitude to the range of answers that I
will accept for this problem since the range of your backgrounds are
relatively diverse.  I am looking for some drawings of the geothermal state.
Starting with a simple cartoon of the geologic process is likely to help you
develop your thoughts and ideas.

You may find that your specific description differs from the hypothetical
surface heat flow detailed in Figure~\ref{fig:thermtec} because these
hypothetical models include a temporal evolution that will likely not be
addressed in your simple models.  Also there are additional processes at
play that during their evolution (especially the thrust) are not accounted
for in most simple models.

The thrust is perhaps the most complicated of the scenarios in Figure
\ref{fig:thermtec} not because of the initial complexities but because the
thermal evolution of a thrust sheet post thrusting typically involves a
number of additional processes.

The original thrusting event creates a duplication of the upper part of
the lithospheric column over one which is undisturbed.  As a result, the
upper part of the geotherm is duplicated resulting in a fir tree shape
to the geotherm, placing hot crust at the base of the thrust
sheet on top of cold crust.  This rapid difference in temperature
changes the fastest with a cold front propagating both upward and hot
front downward from the thrust.  Since the cold front is propagating
upward, it causes the heat flow to initially drop.  After some time, the
heat flow will begin to rise as temperatures within most of the
lithosphere rise above the initial values.
\begin{figure}[hbt]
    \centering
    \includegraphics[]{figures/thrusting.eps}
    \caption{Thrusting and duplication of the crust.}
    \label{fig:thrust}
\end{figure}

If the thermal state eventually returned to the pre-thrust state, then
the heat flow would asymptotically rise to the pre-thrust heat flow.
With heat production within the lithosphere, thickening the crust 
increases in the total radiogenic heat production of the
lithosphere resulting in a higher equilibrium heat flow than the
pre-thrust state.  This explains the long term higher heat flow shown in
Figure~\ref{fig:thermtec}.

But why would heat flow rise so drastically after the initial thrusting?
Assuming the upper crust is more felsic, the upper crust is more
susceptible to melting at relatively modest temperatures.  Rising
temperatures in the formerly upper crust below the thrust can reach
temperatures equivalent to the lower crust and as a result melting.  The
melting is further driven by additional heat created by radiogenic
decay.  The melts will migrate towards the surface advecting heat,
creating high heat flow at the surface.  Because of the additional
radiogenic contribution to heat flow, the decay of heat is slower than
expected from a rift or an arc.

